\chapter{走近大语言模型：能力边界与主案例}

\begin{introduction}
\item {\heiti 本章核心问题}：在企业知识助手这个主案例里，大语言模型究竟是什么，它能做什么、不能做什么，又为什么不能被直接当成一个完整业务系统？
\item {\heiti 主案例推进到哪一步}：定义企业知识助手的第一版目标，明确模型层、任务层、系统层三类能力边界。
\item {\heiti 读完后能做什么}：遇到一个 LLM 问题时，先判断它属于模型能力、任务能力还是系统能力，再决定优先修哪里。
\end{introduction}

\section{学习目标}
\begin{itemize}
  \item 从工程视角理解大语言模型在业务系统中的角色，而不是只把它当成一个热门名词。
  \item 建立模型能力、任务能力和系统能力三层边界，避免把所有问题都推给同一个技术手段。
  \item 理解为什么本书要用“企业知识助手”作为贯穿案例，以及这个案例的第一版应当做成什么样。
  \item 学会一套最小决策框架，在提示工程、RAG、微调、继续预训练和系统治理之间做初步取舍。
\end{itemize}

\section{问题背景}
很多团队第一次接触大语言模型，通常从几个令人惊艳的演示开始：它会写邮件、会总结会议纪要、会解释代码，甚至能和人多轮对话。演示看起来很好。但演示成功不等于系统可用。员工在企业内部助手里提问”出差住宿超标还能不能报销”，系统给出一段流畅的回答，却没引用制度来源；制度更新一周后，它仍沿用旧版本；同一问题换一种问法，又给出另一个答案。到这一步，问题已经不是”模型会不会说话”，而是”系统能不能可信地交付答案”。

初学者进入 LLM 领域时，几乎都会经历这种落差。表面上看，模型输出了自然语言；真正落地时，团队却要面对来源追溯、权限控制、时效更新、错误复盘、延迟成本和版本回归等一整套问题。大语言模型带来的不是一个单点工具，而是一条新的工程链路。这条链路上的每一环出问题，最终都表现为"系统回答不对"。

没有一条清晰主线，学习 LLM 很容易碎片化。有人记住了很多术语，却说不清它们在系统里的位置；有人把提示工程、RAG、微调和对齐训练混成一个”让模型更好”的大筐；也有人过早沉迷某个专题——比如 LoRA 或 GraphRAG——却始终没建立从底层模型到业务交付的整体地图。本章的任务不是罗列名词。它要先搭一个坐标系，让后面每一章都有落点。

\section{核心概念}
\subsection{什么是工程视角下的大语言模型}
从数学上说，大语言模型首先是一个条件概率模型。给定前文 \(x_{<t}\)，它学习下一个 token 的条件分布
\[
p_\theta(x_t \mid x_{<t}).
\]
这句话的工程含义比它看起来大得多。模型最原始的能力，不是”理解世界真相”，而是”根据已有上下文继续产生最可能的序列”。它擅长压缩模式、匹配语境、延续分布。但它不天然拥有外部事实、权限规则和实时状态——这些东西从来不在训练目标里。

\begin{definition}[工程视角下的 LLM]
在工程系统里，大语言模型可以被视为一个可调用的语义计算核心。它负责阅读上下文、生成文本、组织解释、输出结构化结果或完成一定程度的推理，但它本身不天然拥有业务真相、外部状态、权限约束和可追溯记忆。
\end{definition}

这个定义决定了后面所有章节的分工。把模型直接当成数据库、工作流引擎和权限系统的总和，项目几乎一定会在上线后出问题。但如果把模型看成一个强大的语义处理部件，再为它补上知识接入、任务约束和系统治理，路线反而会清晰很多。

\subsection{为什么一次惊艳演示不等于系统能力}
刚接触 LLM 的团队很容易被一两个漂亮样例带偏。模型在演示里答对了差旅制度、周报写作、代码解释，团队下意识就觉得：既然它已经会答，那离上线应该不远了。

这个判断的漏洞在于，它把点状成功误当成了分布上的稳定能力。

企业知识助手真正要面对的不是一道精心挑选的问题，而是一整片真实查询分布。同一个制度会被不同部门用不同口语追问；不同版本文档彼此冲突；部分问题根本没有足够材料；有些问题即使答案大体正确，也必须带引用、带边界、带转人工建议。演示里的成功只说明”模型在这个点上能做对”。它说明不了”系统在这一类问题上能稳定交付”。

\begin{table}[H]
\centering
\small
\begin{tabularx}{\textwidth}{>{\raggedright\arraybackslash}p{2.7cm}>{\raggedright\arraybackslash}p{3.1cm}>{\raggedright\arraybackslash}X}
\toprule
\textbf{演示里看到的现象} & \textbf{为什么容易误导} & \textbf{上线前真正还缺什么} \\
\midrule
单题回答很流畅 & 流畅会被误当成可靠 & 还缺事实校验、来源追溯与错误兜底 \\
改写、总结效果很好 & 易让人以为理解已经够用 & 还缺任务边界、格式约束与回归验证 \\
几轮对话看起来顺畅 & 会被误当成长期记忆稳定 & 还缺上下文裁剪、状态管理和记忆治理 \\
少量示例上答对企业问题 & 会被误当成“懂公司业务” & 还缺知识接入、权限过滤和文档时效更新 \\
\bottomrule
\end{tabularx}
\caption{演示成功回答的是“这个样例”，系统成功回答的是“这一类问题”}
\label{tab:demo-vs-system}
\end{table}

本书后面反复强调主案例、评测和回放，不是在给工程增加负担，而是在补这一层最容易被忽略的鸿沟——会回答，不等于能交付。

\subsection{从“会说话”到“能交付”之间隔着系统设计}
一个会说话的模型和一个能交付的业务系统之间，至少隔着四道门槛：
\begin{itemize}
  \item 事实门槛：模型说的话是否与当前业务事实一致。
  \item 边界门槛：模型是否知道什么时候该回答，什么时候该拒答或转人工。
  \item 来源门槛：结论能不能回到证据、文档或工具输出。
  \item 工程门槛：系统能不能被监控、回放、评测和迭代。
\end{itemize}

初看上去像”模型不够聪明”的问题，多数其实出在后面三道门槛上。企业知识助手之所以适合做贯穿案例，就是因为它会同时暴露这四道门槛——不是停留在”写一段漂亮回复”的层面就够了。

\subsection{知识、规则与实时状态不是同一种东西}
很多团队一开始会把企业里的所有信息都统称为”上下文”，然后默认用同一种方式接给模型。想法很自然，但工程上往往出问题。企业系统里至少同时存在三类完全不同的信息来源：
\begin{itemize}
  \item \textbf{知识}：例如制度正文、产品手册、FAQ、操作说明，特点是可写成文档、可被检索、需要引用来源。
  \item \textbf{规则}：例如审批链必须经过谁、什么金额触发什么校验、哪些字段不能为空，特点是它们更像流程约束，而不是一段“可自由解释的文字”。
  \item \textbf{实时状态}：例如某张单据当前走到哪一步、某个权限是否已开通、某台服务此刻是否健康，特点是它们随时间变化，往往要通过 API、数据库或工具查询获得。
\end{itemize}

模型表面上都能把这三类信息“读成文字”，但系统接入方式不应该一样。知识更适合走检索与引用链，规则更适合落在工作流和校验器里，实时状态则更适合工具调用或系统查询。只要把这三类东西混在一起，团队就很容易把本该由规则系统保证的事情交给 prompt，把本该实时查询的事情塞给参数记忆，最后得到一套看似会说、其实不可靠的助手。

本书后面把 RAG、工具调用、评测和系统工程分开讲，原因就在这里——它们不是在争夺”谁更先进”，而是在分别补三类不同的信息缺口。读者只要先接受这个分法，后面就更容易理解为什么有些问题该修文档接入，有些该修流程约束，还有些必须去查实时系统。

\subsection{企业里最危险的错答，不是答不上来，而是答得像真的}
在通用聊天场景里，模型偶尔答不上来也许只是体验问题；但在企业环境里，更危险的往往是另一种情况：它答得很流畅、很完整、甚至很像有依据，但关键事实、适用范围或版本边界其实是错的。

这种风险之所以大，是因为流畅性会放大信任。员工看到一段条理清楚的回答，很容易默认它已经查过制度；支持团队看到一段”像资深同事写的解释”，也容易忽略其中是否夹带了旧版本流程或缺失的例外条件。系统最先需要学会的能力，往往不是”尽量答更多”，而是”在证据不足时别把猜测包装成结论”。

所以第一章就要把一个看似保守的原则说清楚：拒答、澄清和回指证据，不是模型弱，而是系统可信的起点。后面无论是提示工程、RAG、评测还是对齐训练，都在帮系统把这条底线守稳。

\subsection{主流模型家族与任务差异}
从工程上看，LLM 并不是只有一种统一架构。粗略地说，主流模型可以分成三类：
\begin{itemize}
  \item \textbf{Encoder}：更擅长理解任务，例如分类、匹配、抽取和排序，代表是 BERT 一类模型。
  \item \textbf{Encoder-Decoder}：先完整理解输入，再生成输出，适合翻译、摘要、改写和结构化转换，代表是 T5 一类模型。
  \item \textbf{Decoder-Only}：按照自回归方式从左到右生成 token，天然适合开放式生成、对话和通用指令跟随，是当前大语言模型的主流路线。
\end{itemize}

对工程读者来说，重要的不是死记这些名字，而是知道它们偏好的任务类型不同。任务核心是”判断”和”抽取”的，传统 Encoder 模型在成本和延迟上往往更有优势；任务核心是”生成”和”交互”的，Decoder-Only 更自然。企业知识助手以大语言模型为主，原因很具体：它同时需要解释、澄清、摘要、追问和多轮对话，这些能力更贴近自回归生成范式。

\subsection{Base 模型、聊天模型、Embedding 模型与重排序模型不是一回事}
旧稿里分散提到过 base model、chat model、embedding model、reranker 等名词。对入门读者来说，真正要先建立的是一个更实用的判断：\textbf{“模型”不是单数。} 在同一个企业知识助手里，被我们统称为“LLM 系统”的，往往同时会用到几类完全不同职责的模型。

\begin{table}[H]
\centering
\small
\begin{tabularx}{\textwidth}{>{\raggedright\arraybackslash}p{2.7cm}>{\raggedright\arraybackslash}p{3cm}>{\raggedright\arraybackslash}X}
\toprule
\textbf{模型类型} & \textbf{主要负责什么} & \textbf{为什么不能互相替代} \\
\midrule
Base / Decoder-Only 底座模型 & 开放式生成、解释、改写、多轮对话 & 它擅长生成，但不一定最擅长检索、排序或结构化判断 \\
聊天 / 指令模型 & 在底座上补足对话格式、角色遵循和指令执行 & 它更适合交互，但不等于更适合做向量检索或精排序 \\
Embedding 模型 & 把 query 和文档映射到同一向量空间，支撑召回 & 它解决的是“先找回来”，不是“最后怎么回答” \\
重排序模型 & 在候选集合里判断哪几段最值得被送进 prompt & 它更像精排器，不负责开放式生成 \\
OCR / 文档理解模型 & 把扫描件、PDF、表格和版面先读成结构化证据 & 没有它，很多知识根本进不了后续链路 \\
\bottomrule
\end{tabularx}
\caption{企业知识助手里的“模型”通常至少分成生成、检索、排序和文档理解几类}
\label{tab:model-role-map}
\end{table}

这张表要提前拆掉一个常见误解。很多团队一开始会问：”有没有一个更强的模型，换上去以后检索也更准、表格也更稳、排序也更好？” 现实通常不是这样。生成模型再强，也不自动等于最适合做向量索引；检索器再强，也不自动懂企业回答该先给依据还是先给结论。后面本书把 RAG、评测、文档理解和对齐分开讲，就是因为这些职责本来不应该被压成同一个模型的单一责任。

\subsection{同一个企业助手，往往是一套多模型协作系统}
一旦把上面的角色分开，很多工程现象就会自然清楚。用户看到的是一个“企业知识助手”，后台实际上却可能是一整条协作链：先由文档解析器把 PDF 和表格读出来，再由 embedding 模型召回候选，再由重排序模型精排，再由生成模型组织回答，最后再由规则系统和评测门禁做约束。

这个认识有两个直接收益。一是防止团队把所有效果问题都误推给”聊天模型不够强”。二是帮读者理解为什么后面章节看起来在讲很多不同对象，却都服务同一个主案例——企业知识助手从来不是”一颗更聪明的大脑”，而是一套多角色协作、分工明确的工程系统。

\subsection{什么时候根本不该优先上大语言模型}
旧稿里关于“模型选择指南”的经验，在教材版导论里也应该明确保留。工程团队最常见的高成本错误之一，就是把所有语言任务都先默认成”大模型任务”。更稳的判断通常是：
\begin{itemize}
  \item 任务输出空间很封闭——固定标签分类、是否通过审核、字段抽取、近重复检测——优先怀疑小模型、规则或传统检索是否已经够用。
  \item 任务要求极低延迟、极高吞吐，回答形式本身非常标准化，先用 Encoder、小分类器或模板系统，往往比直接上 LLM 更稳。
  \item 任务真正难点在最新事实、权限状态或跨系统查询，先补知识接入和工具链，而不是先升级参数规模。
\end{itemize}

一句话：不是因为任务和语言有关，就一定要用最大、最通用的模型。企业知识助手看起来是一个 LLM 系统，但它底下常常同时依赖小分类器、检索器、重排序器、权限规则和工作流引擎。真正成熟的工程不是”处处都用大模型”，而是知道大模型应该只负责哪一层。

\begin{table}[H]
\centering
\small
\begin{tabularx}{\textwidth}{>{\raggedright\arraybackslash}p{3cm}>{\raggedright\arraybackslash}p{3cm}>{\raggedright\arraybackslash}X}
\toprule
\textbf{任务形态} & \textbf{更稳的起手式} & \textbf{为什么不必先上大语言模型} \\
\midrule
固定标签分类、是否通过审核 & 小分类器 / Encoder 模型 / 规则 & 输出空间封闭，延迟和成本往往比生成能力更关键 \\
字段抽取、表单结构化 & 抽取模型 + 规则校验 & 任务目标明确，先追求可验证和可回填更划算 \\
近重复检测、语义匹配、排序 & 检索器 / 重排序器 & 重点在比较与排序，不在开放式生成 \\
最新制度问答、需要引用依据 & RAG + LLM & 真正难点在知识接入与证据组织，而不是参数记忆 \\
开放式解释、多轮澄清、复杂改写 & LLM 为主，配合系统约束 & 这类任务才真正需要生成、澄清和风格控制能力 \\
\bottomrule
\end{tabularx}
\caption{“要不要先上 LLM”先看任务形态，而不是先看热度}
\label{tab:first-choice-routing}
\end{table}

\subsection{注意力方向与掩码设计会改变“模型擅长什么”}
旧稿里曾把单双向注意力、不同 Decoder 架构和掩码设计分散开讲。对入门读者来说，更该先抓住一个工程结论：模型能不能看见哪些 token，以及在生成时按什么顺序看见，会直接改变它擅长的任务类型。

Encoder 类模型通常允许当前位置同时看见左右文，因此更适合”读完整体再判断”的任务。做分类、相似度匹配、字段抽取或重排序时，模型需要把整段材料一起看完，再决定某个标签或某个候选是否成立。对这类任务来说，完整读取输入比逐 token 续写输出更重要。

Decoder-Only 模型则依赖因果掩码，只能看见当前位置左边已经出现过的 token。这个限制乍看像“看得更少”，但它恰好与文本生成过程一致：模型像写作者一样，一边读已有上下文，一边继续往后写。因此它天然适合对话、解释、开放式问答和长文本续写。企业知识助手的最终交互界面之所以大多基于这一路线，并不是因为它“最先进”，而是因为它和真实回答过程最一致。

Encoder-Decoder 架构则把“读”和“写”显式拆开：Encoder 负责双向读取输入，Decoder 负责按因果方式生成输出，并通过交叉注意力持续参考输入内容。这种结构在翻译、摘要、改写、结构化转换等任务上很自然，因为这类任务本来就是“先读懂，再改写”。它的价值不在于绝对更强，而在于读写分工更清楚。

架构差异不只是”排行榜差异”，而是信息流差异。任务主要是从完整文本里判断标签或抽字段，全局读入往往更重要；任务主要是一步步继续写、解释和互动，因果掩码更自然；任务是先读懂再改写，读写分离通常更稳。下一章会进一步解释这些信息流在 Transformer 结构里如何实现，这里只需要先把它和任务选择连起来。

\subsection{训练目标决定模型的“天性”}
模型训练目标直接影响它后面最擅长的事情。最常见的两类目标是：
\begin{itemize}
  \item 语言模型目标：根据前文预测下一个 token。它强化的是续写、对话和开放式生成能力。
  \item 掩码恢复或去噪目标：恢复被遮盖或打乱的输入片段。它更强调理解、重建和匹配能力。
\end{itemize}

一个模型并不是越大越”万能”，而是会带着训练目标留下的偏好进入下游场景。擅长自然续写的模型，未必天然适合高精度字段抽取；分类任务上很稳的模型，未必适合长篇开放回答。后面讲提示工程、微调和对齐训练时，其实都是在想办法重新塑造或约束这种”天性”。

\subsection{为什么 Decoder-Only 会成为主流}
过去几年里，Decoder-Only 成为大模型主流，并不只是因为“它能聊天”，还因为它在大规模预训练和通用迁移上更符合当前行业路线：
\begin{itemize}
  \item 它与开放式文本生成的实际使用方式一致，训练目标和部署目标更统一。
  \item 它能直接利用海量无标注语料做自监督学习，不依赖任务标签。
  \item 它在 zero-shot 和 instruction following 场景里，更容易把“语言续写能力”迁移成“按要求回答能力”。
  \item 围绕推理加速、量化、适配器和服务部署的生态最成熟。
\end{itemize}

这并不意味着其他架构没有价值，而是说明：如果一本书的主线是“构建企业知识助手”，那么围绕 Decoder-Only 组织知识最贴近当下的工程现实。

\subsection{长上下文不等于无限记忆}
旧稿在长文本问题上给过很多零散提醒，这里需要在导论中先收成一个清楚判断：上下文窗口变长，只是让模型有机会看更多内容，并不自动意味着它能稳定找到重点、记住重点、引用重点。

工程上至少要同时看三件事：
\begin{itemize}
  \item 能不能放进去：这是窗口长度、显存和延迟问题。
  \item 放进去后会不会被噪声稀释：这是注意力分散和证据密度问题。
  \item 放进去后能不能稳定回指关键依据：这是检索、排序和提示编排问题。
\end{itemize}

长上下文更像一张变大的工作台，而不是一个天然更可靠的记忆系统。只要信息来源是动态知识、跨文档事实或大量制度片段，团队就不该把”窗口加长”误当成”知识接入已经解决”。后面本书把长上下文、RAG、上下文压缩和推理成本放到一条连续主线上讲，道理就在这里。

\subsection{闭卷记忆、上下文记忆与外部知识是三回事}
用户经常会问”模型知不知道这件事”。从工程视角看，这个问题如果不拆开，会直接把团队带进错误路线。LLM 至少有三种完全不同的信息来源：参数里的闭卷记忆、当前上下文里的临时记忆，以及系统额外接入的外部知识。

所谓闭卷记忆，指的是模型在预训练和后训练阶段被参数吸收进去的模式与事实。它能支持常识、写作风格、一般性知识和部分历史事实，但它不是实时更新的，也不带企业权限边界。上下文记忆指的是你在当前 prompt 里实际喂给模型的材料，例如制度条款、日志片段、表格行或对话历史。外部知识则更进一步，意味着系统通过检索、数据库、API 或工具调用，把本来不在当前参数和上下文里的信息临时接进来。

把这三者混在一起，是很多项目早期最昂贵的误会之一。团队误以为“模型知道公司制度”，其实它只是在闭卷记忆里学过一些公开制度写法；团队误以为“对话历史会永久记住”，其实上下文一旦被裁掉，模型就不再持有那段状态；团队误以为“RAG 只是锦上添花”，其实只要答案依赖最新企业知识，它往往就是主路径而不是配角。

一个更成熟的追问方式不是”模型知不知道”，而是：
\begin{itemize}
  \item 这条信息应该来自参数、来自当前上下文，还是来自外部系统？
  \item 如果来自外部系统，证据入口、权限边界和更新节奏是什么？
  \item 如果来自上下文，窗口预算是否真的足以稳定容纳它？
\end{itemize}

这套区分会在全书里反复出现。它解释了为什么“长上下文”不能替代知识接入，也解释了为什么企业知识助手的正确答案很少只靠闭卷记忆获得。

\subsection{涌现能力与工程误读}
旧稿里经常出现“涌现能力”这个词。它通常指模型规模、数据规模和训练预算跨过某些阈值后，一些复杂能力会突然明显增强，比如多步推理、代码生成或复杂指令遵循。这个概念值得知道，但更重要的是知道工程上常见的误读是什么。

\begin{itemize}
  \item 误读一：把所有能力提升都归功于参数变大。实际上，数据质量、训练目标、后训练和系统设计同样关键。
  \item 误读二：把一次惊艳演示当成稳定能力。演示是点状现象，工程系统要看分布、稳定性和回归表现。
  \item 误读三：以为模型足够大后，就不再需要检索、评测和权限系统。事实恰恰相反——模型越强，越需要明确它的职责边界。
\end{itemize}

对工程团队来说，最危险的不是低估模型，而是高估模型。高估意味着把本该由知识库、工作流、权限校验和评测体系承担的责任，错误地压给参数本身。

\subsection{三层能力边界}
为了避免后续章节概念混淆，本书把能力边界固定拆成三层：
\begin{enumerate}
  \item \textbf{模型能力}：阅读上下文、生成文本、遵循指令、完成通用推理。
  \item \textbf{任务能力}：通过提示工程、SFT、PEFT、继续预训练或对齐训练后，在某类任务上表现更稳定。
  \item \textbf{系统能力}：通过检索、工具调用、权限系统、日志、评测和部署策略，把模型嵌入业务流程。
\end{enumerate}

这三层边界不是为了做概念分类，而是为了做技术决策。因为同样是“回答不好”，可能完全属于不同层。员工问的事实在最新制度里，但模型答的是旧版本，这通常优先是系统能力问题；模型明明看到了材料，却总是不按要求给引用，这通常更像任务能力问题；连基本摘要和指令跟随都不稳，才更像模型能力问题。

\begin{table}[H]
\centering
\small
\begin{tabularx}{\textwidth}{>{\raggedright\arraybackslash}p{2.8cm}>{\raggedright\arraybackslash}p{3cm}>{\raggedright\arraybackslash}X}
\toprule
\textbf{能力层} & \textbf{核心职责} & \textbf{常见改进手段} \\
\midrule
模型能力 & 会不会读、会不会写、会不会基本推理 & 换模型、换架构、换基础能力更强的底座 \\
任务能力 & 会不会按当前任务规范稳定输出 & 提示工程、SFT、PEFT、对齐训练 \\
系统能力 & 会不会接入知识、工具、权限和监控 & RAG、工具调用、权限过滤、日志与评测 \\
\bottomrule
\end{tabularx}
\caption{企业知识助手中的三层能力边界}
\label{tab:llm-three-layers}
\end{table}

\begin{note}
把“会回答企业文档里的问题”简单归因于模型本身，通常是错误的。很多时候，这种能力来自检索系统、文档解析、权限过滤和提示编排，而不是参数记忆。
\end{note}

\subsection{Agent 不是第四层能力，而是系统层的组织方式}
很多初学者读到 Agent、工作流编排、工具调用时，会下意识把它理解成”比模型、任务、系统更高级的一层”。这种理解会制造混乱。更稳的说法是：Agent 不是第四层能力，而是系统层把模型、工具、记忆和决策规则组织起来的一种形态。

一个企业知识助手如果带有 Agent 行为，通常意味着它不只生成一段答案，还会做几件额外的事：决定要不要检索、决定要不要调用工具、决定要不要追问用户、决定什么时候升级人工。这些动作听起来很”聪明”，但它们仍然要落回前面三层边界：
\begin{itemize}
  \item 调不调用工具，首先是系统层设计；
  \item 工具结果如何组织进回答，是任务层约束；
  \item 模型能否读懂工具返回值并组织语言，才是模型层能力。
\end{itemize}

很多 Agent 失败案例，真正问题不在”Agent 理念不对”，而在系统层没有把可调用工具、权限边界、失败回退和日志复盘设计清楚。先把这个边界捋顺，后面读到工具调用、RAG 编排和工作流时，读者就不容易误以为自己又进入了一个全新的技术宇宙。

\subsection{先判断问题属于哪一层}
对初学者来说，最容易犯的错不是”不知道术语”，而是不知道当前问题属于哪一层。一套更稳妥的判断顺序是：
\begin{enumerate}
  \item 答案依赖最新制度、权限边界、外部状态或跨系统结果，优先当成系统能力问题。
  \item 材料通常是够的，但输出格式不稳、语气不合规、引用规则执行不好，优先当成任务能力问题。
  \item 连通用阅读、摘要和指令跟随都表现很差，再优先怀疑模型能力或选型本身。
\end{enumerate}

这套顺序看起来简单，却会直接决定项目成本。很多团队本可以先修检索和知识接入，却误判成“模型太弱”；也有不少团队本该先统一输出规范，却一上来就做复杂训练。

\begin{table}[H]
\centering
\small
\begin{tabularx}{\textwidth}{>{\raggedright\arraybackslash}p{3.2cm}>{\raggedright\arraybackslash}p{2.8cm}>{\raggedright\arraybackslash}X}
\toprule
\textbf{观察到的现象} & \textbf{优先归层} & \textbf{第一动作} \\
\midrule
引用过期制度或错误知识源 & 系统层 & 先查知识源更新、检索链路和权限过滤 \\
材料充足但总不按格式答题 & 任务层 & 先查提示模板、SFT 样本和输出约束 \\
连基础摘要和跟随指令都不稳 & 模型层 & 先查模型选型或基础能力是否够用 \\
同一问题换说法就明显失效 & 系统层或任务层 & 先查检索表达对齐，再查提示约束 \\
\bottomrule
\end{tabularx}
\caption{从现象反推能力层级的最小判断表}
\label{tab:llm-routing}
\end{table}

\begin{figure}[H]
\centering
\begin{tikzpicture}[
  font=\small,
  box/.style={llm box, minimum width=3.4cm, minimum height=1.2cm},
  arr/.style={llm arr}
]
\node[box] (model) {模型能力\\读上下文\\生成文本};
\node[box, right=1.2cm of model] (task) {任务能力\\提示工程\\微调与对齐};
\node[box, right=1.2cm of task] (system) {系统能力\\检索、权限、日志\\评测与监控};
\draw[arr] (model) -- (task);
\draw[arr] (task) -- (system);
\node[llm note box, below=0.9cm of task, minimum width=10.4cm, minimum height=1.2cm] (assistant) {企业知识助手第一版目标：高频问题有依据地回答，依据不足时明确拒答，且整条链路可以被复盘};
\draw[arr] (task) -- (assistant);
\end{tikzpicture}
\caption{企业知识助手的三层能力边界}
\label{fig:assistant-boundary}
\end{figure}

\subsection{误判问题层级，会把团队带进错误的高成本路线}
三层能力边界最重要的价值，不只是帮助读者“分清概念”，而是避免团队在错误方向上投入过多成本。因为同样一句“系统回答不够好”，如果归层归错，后续动作往往会完全跑偏。

最典型的三类误判是：把知识时效问题误诊成模型太弱，于是急着做微调；把输出格式和引用顺序问题误诊成检索不够强，于是越堆越多上下文；把工具失败和权限校验缺失误诊成 prompt 没写好，于是在系统层缺口没补上的情况下不断叠规则。表面上看，团队一直在“优化模型”；实际上，它只是绕着真正问题在打转。

\begin{table}[H]
\centering
\small
\begin{tabularx}{\textwidth}{>{\raggedright\arraybackslash}p{3.1cm}>{\raggedright\arraybackslash}p{3.3cm}>{\raggedright\arraybackslash}X}
\toprule
\textbf{真实问题} & \textbf{常见误诊} & \textbf{被带偏后的高成本后果} \\
\midrule
制度更新后仍答旧内容 & 误诊成底座模型太弱 & 花大量时间微调，却仍然答旧知识 \\
回答顺序乱、引用总在最后 & 误诊成检索证据不够多 & 上下文越塞越长，真正的格式问题反而更糊 \\
工具查询失败或权限没校验 & 误诊成提示词没写细 & 规则叠得越来越长，但越权风险并未下降 \\
同类任务长期不稳定 & 误诊成“再换个更大模型” & 成本升高，却没有建立任务层约束与评测闭环 \\
\bottomrule
\end{tabularx}
\caption{归层错误最大的代价，不是概念混乱，而是优化路线被带偏}
\label{tab:misdiagnosis-cost}
\end{table}

本章把”先判层，再选手段”放在全书最前面，原因就在这里。对工程入门读者来说，学会停止在错误层级上过度投资，本身就是一种关键能力。

\subsection{什么问题更适合提示、RAG、微调或继续预训练}
旧稿中大量内容彼此交叉，本书需要在第一章就先给出一个最小判断框架，否则后面每一章都会显得像”另一个让模型变强的方法”。在展开之前，先用一句话说清这几个术语：

\begin{itemize}
  \item \textbf{提示工程}（Prompt Engineering）：不改模型参数，只通过修改输入文本（系统提示、示例、格式要求）来控制模型行为。
  \item \textbf{RAG}（Retrieval-Augmented Generation，检索增强生成）：先从外部知识库中检索相关文档片段，再把它们塞进模型的输入上下文，让模型基于真实材料回答。详见第 7 章。
  \item \textbf{微调}（Fine-Tuning，也叫 SFT——Supervised Fine-Tuning，监督微调）：用一批标注好的”输入→期望输出”样本，继续训练模型参数，让它在特定任务上表现更稳定。详见第 5 章。
  \item \textbf{继续预训练}（Continued Pre-Training）：用大量领域文本（不需要标注）继续训练模型，让它熟悉特定领域的语言分布和术语。详见第 6 章。
\end{itemize}

有了这些定义，下面的判断框架就更容易理解：
\begin{itemize}
  \item \textbf{先用提示工程}：当问题主要是输出格式、角色约束、风格、步骤展开或短上下文内的规则执行。
  \item \textbf{优先用 RAG}：当问题主要依赖最新文档、企业知识、表格内容、外部状态和可追溯证据。
  \item \textbf{考虑微调}：当同一类任务大量重复出现，且提示已经很长、很脆弱、很难维护。
  \item \textbf{考虑继续预训练}：当模型对某个领域语言、术语、文风和长文档分布明显陌生。
\end{itemize}

后面我们会分别展开这些手段，但在第一章先记住一件事就够了：它们不是彼此替代的流行选项，而是在解决不同层级的问题。

\subsection{质量、延迟、成本与风险不会同时最优}
初学者最容易写出的需求，往往也是最不可能同时满足的需求：“回答要更准、更快、更便宜、覆盖更广，而且风险还要更低。” 从工程角度看，这四件事几乎从来不会同时往同一个方向走。只要系统开始落地，团队就必须学会在质量、延迟、成本和风险之间做取舍。

例如，把更多文档塞进上下文，可能提高覆盖，却会增加延迟和成本；把拒答阈值调得更保守，可能降低风险，却会牺牲部分表面上的“回答率”；换更大的底座模型，也许能提升表达质量，却会立刻抬高推理预算与部署复杂度；引入更多工具和工作流，也许能让答案更接近真实业务，却会让链路更长、更难回放。

这并不意味着系统只能”顾此失彼”，而是意味着第一版必须先声明谁是主目标，谁是约束条件，谁可以暂时让步。对企业知识助手来说，第一版通常更应该优先守住”有依据、可复盘、风险可控”，而不是急着在所有长尾问题上追求看起来很聪明的覆盖率。

当团队讨论”为什么不直接上更强模型””为什么不把所有知识都塞进去”时，更成熟的回答通常不是”技术上做不到”，而是”这样会把哪个代价推高到当前阶段不能接受”。工程教材和产品演示最大的区别就在这里：教材要教会读者看到取舍，而不是只看到能力。

\subsection{专业领域不一定先做“行业大模型”}
旧稿里还反复讨论过一个容易引发误判的问题：做企业应用，是不是就要先做一套自己的行业大模型。教材版里更适合把结论说得直接一些：大多数团队的第一步都不该是自造基础模型，而该是先判断问题到底属于知识接入、任务规范，还是语言分布。

只有在下面几类情况下，继续预训练、专属底座甚至更重的模型改造才会开始合理：
\begin{itemize}
  \item 领域术语、文档格式、日志模式与通用语料差异极大，模型连文本都读不顺；
  \item 业务规模足够大，长期成本足以覆盖模型治理、训练和评测投入；
  \item 安全隔离、私有化要求或部署环境，使通用闭源服务难以满足要求。
\end{itemize}

反过来，如果主要问题是“制度常变”“知识要可追溯”“回答要带引用”“权限边界要保守”，那优先级通常仍然是 RAG、工具调用、评测和工作流治理，而不是先造一套更大的参数系统。对工程入门读者来说，学会不被“行业大模型”这个词带偏，本身就是第一章应该完成的任务之一。

\subsection{为什么必须有一个贯穿案例}
一本工程教材如果缺少持续推进的业务对象，知识就会碎片化。本书固定使用”企业知识助手”作为主案例。它的目标不是做一个”无所不能”的聊天机器人，而是做一个可上线、可治理、可迭代的知识系统。

这个案例合适，是因为它天然会触发整本书的核心主题：前半书解释模型为什么会回答、怎样控制回答；中段解释如何接入企业知识；后段解释怎样评测、对齐、部署和扩展。它能把”理解模型”到”交付系统”这一整条学习路径串起来。

\section{主案例推进}
这一章里，我们把企业知识助手的第一版目标故意收得很窄。第一版不是要覆盖全公司所有知识，不是要做复杂代理系统，也不是要替代人工审批。它只需要完成三件事：
\begin{itemize}
  \item 对高频制度和产品问题给出可读答案。
  \item 当依据不足时明确说“不足以判断”。
  \item 让团队能够通过日志和评测知道它为什么答对、为什么答错。
\end{itemize}

\subsection{第一版的验收标准不是更聪明，而是更可控}
很多项目在第一版就忍不住追求”覆盖更多场景””回答更像专家””最好还能自动行动”。从工程角度看，这恰恰是最需要克制的地方。第一版系统最值得追求的，不是它显得多聪明，而是它的行为是否够稳定、够保守、够可复盘。

\begin{table}[H]
\centering
\small
\begin{tabularx}{\textwidth}{>{\raggedright\arraybackslash}p{2.8cm}>{\raggedright\arraybackslash}p{3.2cm}>{\raggedright\arraybackslash}X}
\toprule
\textbf{验收维度} & \textbf{第一版最小要求} & \textbf{为什么先只做到这里} \\
\midrule
回答可读性 & 高频问题能给出清晰、短而完整的回答 & 先让系统能被真实用户使用，而不是只会演示 \\
依据边界 & 材料不足时能明确说“不足以判断” & 先压住高风险误答，而不是盲目追求覆盖率 \\
来源追溯 & 至少能说明依据来自哪类文档或系统 & 为后续 RAG、评测和审计打基础 \\
回放能力 & 出错样例能被日志和固定输入重现 & 没有复盘，就谈不上系统迭代 \\
角色边界 & 不直接替代审批、付款或权限发放 & 先把风险责任收住，避免“会答就敢做” \\
\bottomrule
\end{tabularx}
\caption{企业知识助手第一版先验收“可控”，再验收“更强”}
\label{tab:v1-acceptance}
\end{table}

这张表和后面第 9 章的正式评测体系并不重复。这里先给读者一个导论级的判断：第一版是否值得进入下一章，不是看它像不像一个聪明聊天机器人，而是看它有没有最小的系统纪律。

\subsection{先定义用户、问题和风险}
为了让案例后续可推进，我们先把第一版的用户和问题收敛出来。目标用户主要是员工和内部支持团队，问题聚焦在制度说明、权限流程、产品 FAQ 和运维常见问答。之所以这么做，不是因为这些场景最容易，而是因为它们同时满足三个条件：高频、可验证、上线价值明确。

\begin{table}[H]
\centering
\small
\begin{tabularx}{\textwidth}{>{\raggedright\arraybackslash}p{2.8cm}>{\raggedright\arraybackslash}p{3.2cm}>{\raggedright\arraybackslash}X}
\toprule
\textbf{维度} & \textbf{第一版选择} & \textbf{原因} \\
\midrule
目标用户 & 员工、支持团队、运维同学 & 问题高频，反馈快，价值明确 \\
问题范围 & 制度、FAQ、流程、手册 & 有文档依据，便于检索、引用和评测 \\
交付目标 & 回答、引用、拒答、可复盘 & 先把可信度做稳，再谈复杂自治能力 \\
\bottomrule
\end{tabularx}
\caption{企业知识助手第一版的案例边界}
\label{tab:assistant-v1-scope}
\end{table}

\subsection{先选权威知识源，不把“能搜到的文本”都接进来}
企业知识助手第一版还有一个很容易被忽略、但决定系统上限的动作：先选知识源，再做模型链路。如果这一步不收紧，后面再好的检索、提示和微调，也只是在混杂材料里做更复杂的猜测。

对第一版来说，更稳的知识源通常只有几类：正式制度库、经过版本管理的产品 FAQ、可追溯的流程文档、带负责人和更新时间的运维手册。相反，群聊截图、个人笔记、会议口头结论、过期培训材料和无法确认版本的二手整理，默认都不应直接作为主知识源。不是说它们永远没用，而是第一版系统没有必要先把最脏、最难追责的材料全吃进来。

这个选择和后面 RAG 章节并不重复。RAG 解决的是“怎么接入知识”，而这里先回答“哪些材料有资格被接入”。对教材读者来说，这一步非常重要，因为很多所谓“模型幻觉”其实并不是模型凭空乱说，而是系统把互相冲突、版本不明或责任主体不清的文本一起喂了进去。

因此，第一版的正确顺序通常是：先锁定少量权威知识源，再建立更新节奏和失效机制，最后才讨论怎么把这些知识稳定接给模型。这样做看起来像是在缩小系统能力，实际上是在给后面所有章节争取一个干净可治理的底座。

\subsection{第一版明确不做什么}
一开始就明确“不做什么”，和明确“要做什么”一样重要。企业知识助手第一版默认不做以下事情：
\begin{itemize}
  \item 不直接替代审批、报销或权限发放流程。
  \item 不在证据不足时强行给确定结论。
  \item 不把历史对话当成永久知识。
  \item 不试图覆盖所有部门的全部长尾知识。
\end{itemize}

这些限制不是保守，而是工程上的收敛。教材里很多章节之所以会显得杂乱，就是因为没有先定义边界，导致任何技术都能被说成“未来也许有用”。但真正能交付的系统，恰恰是靠边界活下来的。

\subsection{为什么第一版先做高频低风险，而不是最难的问题}
很多团队做内部助手时，最容易被“最难的问题”吸走注意力。例如一上来就想解决跨部门复杂审批推理、几十页制度综合判断、历史多轮对话长期记忆，或者直接让系统自动触发动作。它们当然重要，但如果把这些场景拿来定义第一版，项目很可能会在还没建立评测、回放和边界控制之前，就被复杂度压垮。

第一版优先选择高频低风险问题，原因不只是”更容易做”，而是它同时满足三条工程条件：
\begin{itemize}
  \item 反馈快：高频问题每天都会出现，团队能很快积累真实日志和坏例子。
  \item 答案可验证：制度、FAQ 和手册型问题通常能回到文档依据，适合建立最早的一批评测集。
  \item 风险可控：即使答错，多数场景也可以通过拒答、引用和转人工把损失控制住。
\end{itemize}

反过来看，那些”最难也最酷”的场景，往往同时带着低频、难验证和高风险三重特征。不是不值得做，而是不应该在系统还没有最小治理能力时先做。把第一版范围收窄，不是在压低目标，而是在先把一套可学、可测、可迭代的工程节奏搭起来。

\subsection{本书后续会如何推进这个案例}
从这一章开始，后面的每一部分都会推进同一个系统：
\begin{enumerate}
  \item 先理解模型的语言生成机制，知道它为什么能读上下文并继续回答。
  \item 再学会调用现成模型，完成最小推理与调试闭环。
  \item 接着通过提示工程、微调和继续预训练，让模型更贴近企业任务。
  \item 然后引入 RAG 和文档理解，把企业知识稳定接入系统。
  \item 之后建立评测、对齐和性能优化机制，让系统进入可维护状态。
  \item 最后补齐推理部署和大规模训练工程视角，使系统能够持续扩展。
\end{enumerate}

如果换成更工程化的说法，第一版并不追求“最聪明”，而是追求三个约束同时成立：
\begin{enumerate}
  \item \textbf{可追溯}：答案能回到具体材料，而不是只凭参数记忆。
  \item \textbf{可控}：面对高风险或依据不足问题时，系统知道何时拒答、何时转人工。
  \item \textbf{可复盘}：回答好坏都能被日志、评测和样本回放解释。
\end{enumerate}

\section{常见坑与工程取舍}
\begin{itemize}
  \item \textbf{误把模型当数据库}：模型参数不是实时知识源，业务真相不能只靠预训练记忆。
  \item \textbf{误把 RAG 当万能药}：检索可以补知识，但不能自动解决权限、工作流和输出规范问题。
  \item \textbf{误把微调当前置步骤}：很多场景先用好提示词和检索，收益更高、成本更低。
  \item \textbf{误把单次演示当系统能力}：一次成功回答不代表系统具备稳定交付能力，必须有评测和监控。
\end{itemize}

再补三类在旧稿里反复出现的误区：
\begin{itemize}
  \item \textbf{把模型家族差异理解成排行榜差异}：选型不是挑最火模型，而是看任务类型、成本预算和部署边界。
  \item \textbf{把生成能力等同于理解能力}：能写长答案的模型，不一定适合高精度分类和抽取。
  \item \textbf{把长上下文能力等同于“什么都能塞进去”}：上下文越长，越考验输入质量、检索质量和预算控制。
\end{itemize}

\section{本章练习}
\begin{exercise}[能力分层判断]
某团队的企业助手经常引用过期制度，但回答口吻和格式都很稳定。这个问题首先应归入哪一层能力问题？最优先的修复方向是什么？
\end{exercise}

\begin{solution}
这首先是{\heiti 系统能力问题}，因为错误来自知识时效而不是表达风格。最优先的修复方向是补检索链路、更新知识源和建立依据引用，而不是先做微调。
\end{solution}

\begin{exercise}[先判层再判手段]
某团队发现模型回答时经常漏掉“依据不足时请明确说明无法判断”这条规则，但给定材料通常是充足的。此时应先改什么？为什么不该直接上 RAG？
\end{exercise}

\begin{solution}
应先把问题当成{\heiti 任务能力问题}，优先修改提示模板或补一小批 SFT 样本，强化“依据不足则拒答”的行为。因为当前材料通常已经充足，问题不在知识接入，而在模型没有稳定执行任务规则。
\end{solution}

\section{延伸阅读}
\begin{itemize}
  \item \textbf{Stephen Wolfram, \emph{What Is ChatGPT Doing ... and Why Does It Work?}}：适合理解“下一个 token 预测”和语言行为之间的关系。
  \item \textbf{OpenAI、Anthropic、Google 等模型系统卡}：适合理解能力边界、安全边界和部署边界如何被写成工程约束。
  \item \textbf{Chip Huyen, \emph{AI Engineering}}：适合理解模型、数据与系统能力为什么必须一起设计。
  \item \textbf{企业知识治理与文档生命周期资料}：适合理解为什么很多“模型错误”其实源于知识治理不完整。
\end{itemize}

\section{小结}
本章建立了全书最重要的坐标系：大语言模型是语义计算核心，但不是完整系统；系统能力来自模型、任务约束、知识接入和工程治理的协同；企业知识助手将作为贯穿案例，持续承载这些能力的落地过程。只要先学会判层，再学会选手段，后面的每一章就不再是零散技巧，而会落回同一条工程主线。

\section{下一章过渡}
既然企业知识助手的底层仍然是一个语言模型，那么下一步就必须回答：模型究竟如何“读上下文并继续生成”？下一章将聚焦 Transformer 的核心机制，用最少但足够的原理为后面所有工程章节打基础。
